\documentclass[10pt]{article}
\usepackage[T1]{fontenc}
\usepackage[utf8]{inputenc}
\usepackage[french]{babel}
\usepackage{amsmath,amssymb,mathrsfs,mathpazo}
\usepackage{geometry}
\geometry{a4paper,margin=2cm}
\usepackage{graphicx,array,tabularx,tabularray,colortbl,multirow,multicol,booktabs,enumitem,pifont,ulem,fancybox,diagbox,cancel}
\usepackage{pstricks,pstricks-add,pst-plot,pst-tree,pst-node,pst-text,pst-func,pst-math,pst-bspline,pst-3dplot,pst-eucl,pst-solides3d}
\let\pscalculate\undefined
\usepackage{pst-bezier}
\usepackage{tikz,pgfplots,tkz-tab}
\usetikzlibrary{arrows,arrows.meta,patterns,calc,positioning,trees,decorations.markings,shapes.geometric,shapes.misc,shapes.arrows,intersections,decorations.pathmorphing,decorations.pathreplacing,backgrounds,shadows,babel}
\pgfplotsset{compat=1.18}
\usepgfplotslibrary{fillbetween}
\usepackage[np]{numprint}
\usepackage{listings,makecell,pas-tableur,twoopt}
\let\barre\undefined
\usepackage[pstricks]{tablvar}
\definecolor{darkgreen}{RGB}{0,100,0}
\providecommand{\red}{\color{red}}
\providecommand{\blue}{\color{blue}}
\providecommand{\green}{\color{green}}
\providecommand{\black}{\color{black}}
\providecommand{\white}{\color{white}}
\providecommand{\gray}{\color{gray}}
\providecommand{\cyan}{\color{cyan}}
\providecommand{\magenta}{\color{magenta}}
\providecommand{\yellow}{\color{yellow}}
\providecommand{\violet}{\color{violet}}
\providecommand{\orange}{\color{orange}}
\providecommand{\euro}{\texteuro}
\providecommand{\EUR}{\texteuro}
\providecommand{\ieme}{\textsuperscript{e}}
\providecommand{\degres}{\ensuremath{^\circ}}
\providecommand{\degree}{\ensuremath{^\circ}}
\providecommand{\vv}[1]{\overrightarrow{#1}}
\providecommand{\Coord}[1]{#1}
\providecommand{\decosix}{}
\providecommand{\starredbullet}{\ensuremath{\bullet}}
\pagestyle{empty}
\setlength{\parindent}{0pt}
\usepackage{adjustbox}\DeclareUnicodeCharacter{2236}{:}\DeclareUnicodeCharacter{2190}{\ensuremath{\leftarrow}}
\providecommand{\R}{\mathbb{R}}
\providecommand{\N}{\mathbb{N}}
\providecommand{\Z}{\mathbb{Z}}
\providecommand{\Q}{\mathbb{Q}}
\providecommand{\D}{\mathbb{D}}
\providecommand{\C}{\mathbb{C}}
\providecommand{\e}{{\mathrm{e}}}
\providecommand{\vect}[1]{\overrightarrow{#1}}
\providecommand{\vectt}[1]{\overrightarrow{\mathrm{#1}}}
\providecommand{\vvt}[1]{\overrightarrow{\mathrm{#1}}}
\providecommand{\np}[2][]{{#2~#1}}
\providecommand{\vv}[1]{{\overrightarrow{#1}}}
\renewcommand{\d}{{\mathrm{d}}}
\providecommand{\degres}{{^\circ}}
\providecommand{\degre}{{^\circ}}
\providecommand{\blue}{\color{blue}}
\providecommand{\red}{\color{red}}
\providecommand{\enskip}{\quad}
\providecommand{\llbracket}{\lbrack\!\lbrack}
\providecommand{\rrbracket}{\rbrack\!\rbrack}
\providecommand{\euro}{{\text{€}}}
\providecommand{\no}{{\text{n°}}}
\providecommand{\interval}[2]{{[#1;#2]}}
\providecommand{\intervFO}[2]{{[#1;#2[}}
\providecommand{\dotfill}{\ldots}

\providecommand{\ds}{\displaystyle}
\providecommand{\textding}[1]{#1}
\providecommand{\Oij}{\(\mathrm{O};\vec{\imath},\vec{\jmath}\)}
\providecommand{\Oijk}{\(\mathrm{O};\vec{\imath},\vec{\jmath},\vec{k}\)}
\providecommand{\Ouv}{\(\mathrm{O};\vec{u},\vec{v}\)}
\providecommand{\intervFF}[2]{\ensuremath{[#1;#2]}}
\renewcommand{\euro}{\ifmmode\text{\texteuro}\else\texteuro\fi}
\DeclareUnicodeCharacter{00B0}{\ifmmode^{\circ}\else\textdegree\fi}
\DeclareUnicodeCharacter{20AC}{\ifmmode\text{\texteuro}\else\texteuro\fi}
% Énoncés : Éducation nationale. Transcription et corrigés : APMEP, Denis Vergès et contributeurs. Découpage et classement : CoopMaths, Sébastien Lozano, Éric Elter et contributeurs. Adaptation de la présentation : Maths Explorer.
% CC BY-SA : https://coopmaths.fr/www/about/
% https://coopmaths.fr/alea/static/e3c/2021/tex/e3c_2021_06_ameriquenord_1_cor.tex
\begin{document}
\noindent\textit{APMEP / CoopMaths. CC BY-SA. Présentation adaptée par Maths Explorer.}\par\medskip
\medskip












\begin{tabularx}{\linewidth}{|c|X|X|X|X|}\hline
\textbf{1.} &\multicolumn{4}{|l|}{Pour tout réel $x$,\, $\text{e}^{2x} + \text{e}^{4x}$ est égal à}
\\ \hline
 &\textbf{a.} $\text{e}^{6x}$ &\textbf{b.} $\text{e}^{2x}\left(1 + \text{e}^{2}\right)$ &\textbf{c.} $\text{e}^{3x}\left(\text{e}^{x} + \text{e}^{-x} \right)$ &\textbf{d.} $\text{e}^{8x^2}$
 \\ \hline
\textbf{2.} &\multicolumn{4}{l|}{Dans le plan muni d'un repère orthonormé \Oij, on considère les }\\
& \multicolumn{4}{l|}{vecteurs $\vect{u}(-5~;~2)$ et $\vect{v}(4~;~10)$  et la droite $(d)$ d'équation : $5x + 2y + 3 = 0$.}
\\ \hline
&\textbf{a.} $\vect{u}$ et $\vect{v}$ sont\newline colinéaires&\textbf{b.}  $\vect{u}$ est un\newline  vecteur normal à la droite $(d)$ &\textbf{c.} $\vect{u}$ et $\vect{v}$ sont\newline   orthogonaux&\textbf{d.} $\vect{u}$ est un\newline  vecteur directeur de $(d)$\\ 
\hline
\textbf{3.} & \multicolumn{4}{p{8cm}|}{La dérivée $f'$ de la fonction $f$ définie sur $\R$ par $f(x) = (2x - 1)\e^{-x}$ est :} \\ \hline
&\textbf{a.} \small $2x\e^{-x}$ &\textbf{b.} \small $- 2\e^{-x}$ &\textbf{c.} \small $(- 2x + 3)\e^{- x}$ &\textbf{d.} \small $2\e^{-x} + (2x - 1)\e^{-x}$\\ 
\hline
\textbf{4.} &\multicolumn{4}{l|}{Pour tout réel $x$, on a $\sin(\pi + x) =$}\\ \hline
&\textbf{a.} $-\sin (x)$&\textbf{b.} $\cos (x)$ &\textbf{c.} $ \sin (x)$&\textbf{d.} $-\cos (x)$
\\ \hline
\begin{pspicture*}(0,0)(0.4,5.25) \uput[r](-0.1,4.8){\textbf{5.}} \end{pspicture*}
&\multicolumn{2}{p{5cm}|}{
\vspace{-4cm}Soit $f$ une fonction définie et dérivable sur $\R$ dont la courbe représentative est donnée ci-contre.
La tangente à la courbe au point A est la droite $T$.}&
\multicolumn{2}{|l|}{
\psset{unit=0.75cm}
\begin{pspicture*}(-1,-2.5)(6,4.5)
\psgrid[subgriddiv=1,gridcolor=lightgray](-1,-3)(6,4)
\psaxes[linewidth=1.25pt,labelFontSize=\scriptstyle]{->}(0,0)(-0.99,-3)(6,4)
\psplot[plotpoints=500]{-0.2}{1.5}{3 x 5 mul sub}
\psplot[plotpoints=2000,linewidth=1.25pt,linecolor=blue]{-0.2}{6}{0.9 x  sub 3.3333 mul 2.71828  x 0.45 mul  exp div}
\uput[l](1,-2){$(T)$} \uput[u](3.5,-1.9){$\blue \mathcal{C}_f$} 
\end{pspicture*}
}
\\ \hline
&\textbf{a.} $f'(0) = 3$ &\textbf{b.} $f'(0) =\dfrac{1}{5}$\rule[-12pt]{0pt}{30pt} &\textbf{c.} $f'(0) = 5$ &\textbf{d.} 
$f'(0) = - 5$\\ \hline
\end{tabularx}

\bigskip

\textbf{Solutions}

\smallskip

\textbf{1. c.}

$\e^{3x}\left (\e^{x}+\e^{-x}\right ) = \e^{4x}+\e^{2x}$

\medskip

\textbf{2. c.}

$\vect{u}\cdot\vect{v} = -4\times 5 + 2\times 10 = -20+20=0$ donc $\vect u \perp \vect v$

\medskip

\textbf{3. c.}

$f(x) = (2x - 1)\e^{-x}$ donc $f'(x)= 2\times \e^{-x} + (2x-1) \times (-1)\e^{-x} = (2 - 2x +1)\e^{-x} = (-2x+3)\e^{-x}$

\medskip

\textbf{4. a.}

Les points correspondant à $x$ et $\pi+x$ sur le cercle trigonométrique sont diamétralement opposés  donc $\sin(\pi+x) = - \sin(x)$.

\medskip

\textbf{5. d.}

$f'(0)$ est le coefficient directeur de la tangente $T$ qui passe par les points de coordonnées $(0~,~3)$ et $(1~,~-2)$; donc $f'(0) = \dfrac{-2-3}{1-0}=-5$.

\bigskip

\end{document}
